\documentclass[11pt,a4paper]{article}

\usepackage[margin=1in]{geometry}
\usepackage{amsmath, amssymb}
\usepackage{graphicx}
\usepackage{booktabs}
\usepackage{hyperref}
\usepackage{caption}
\usepackage{subcaption}
\usepackage{float}
\usepackage{microtype}
\usepackage{times}
\usepackage{tabularx}

\hypersetup{
    colorlinks=true,
    linkcolor=blue,
    citecolor=blue,
    urlcolor=blue
}

\title{\textbf{Persuading the Machine: Empathy-Based Appeals, Economic Rationality, and the AI Trust Paradox in Human--Large Language Model Bargaining}}

\author{
    \textbf{Team 4 -- Neuro-Project Research Study}\\
    Department of Computer Science \& Cognitive Science\\
    \texttt{\{research-team\}@university.edu}
}
\date{November 2026}

\begin{document}

\maketitle

\begin{abstract}
As large language models (LLMs) transition from passive informational retrieval tools to autonomous commercial and transactional agents, understanding the psychological and communicative levers that govern their economic behavior is critical. This paper presents an experimental investigation ($N = 101$ participants, 88 completed agreements) evaluating whether empathy-based interpersonal appeals or market-based rational arguments elicit greater financial concessions from an LLM seller. Participants engaged in live, multi-turn negotiations against ``Max'', an autonomous seller powered by Llama-3 (\texttt{openai/gpt-oss-120b} via Groq) governed by an emergency reservation floor of \text{₹}23,000 and an initial listing price of \text{₹}30,000. 

Primary outcome analysis revealed that buyers employing empathy-based appeals secured significantly lower final agreement prices ($M = \text{₹}25,534.56, SD = 988.27$) than those using economic arguments ($M = \text{₹}26,059.73, SD = 884.13$), achieving an average discount advantage of \text{₹}525.18 ($t(83.9) = -2.623, p = 0.0087$, Cohen's $d = -0.561$). Non-parametric Mann--Whitney tests confirmed extreme statistical significance ($U = 562.5, p = 0.0001$). A Baron and Kenny mediation analysis demonstrated that this discount is mediated by elevated perceptions of social warmth (Sobel $z = -2.48, p = 0.013$). Furthermore, multiple regression modeling revealed an unexpected ``AI Trust Paradox'': baseline trust in AI was a powerful predictor of price discounts ($\beta = -553.19, p = 0.0302$), wherein high-trust buyers engaged in reciprocal, incremental bargaining that triggered the LLM's prosocial alignment, whereas low-trust buyers relied on combative ultimatums that activated the agent's programmatic defensive boundaries. These findings demonstrate that conversational vulnerability and cooperative engagement effectively exploit the RLHF agreeableness bias inherent in modern instruction-tuned models.
\end{abstract}

\section{Introduction}
Negotiation represents a quintessential domain of human social cognition, traditionally framed through the dual lenses of economic game theory (focusing on outside options, reservation values, and analytical surplus division) and behavioral psychology (focusing on affect, interpersonal rapport, and perspective-taking). However, the deployment of Large Language Models (LLMs) as autonomous bargaining counterparties creates an unprecedented paradigm: algorithmic negotiation.

A foundational question in human-AI interaction is whether conversational agents respond more readily to human emotional vulnerability or structured market logic. Classical economic theory would suggest that an intelligent agent should discount an asset exclusively when presented with verifiable market comparables, hardware depreciation, and alternative purchase options (BATNA). Conversely, human-centered alignment research reveals that LLMs trained via Reinforcement Learning from Human Feedback (RLHF) exhibit a systematic bias toward politeness, accommodation, and conflict avoidance when interacting with vulnerable human interlocutors.

This study provides an exhaustive empirical investigation ($N = 101$) of this dichotomy, testing:
\begin{enumerate}
    \item Whether empathy appeals yield superior financial payoffs compared to economic arguments.
    \item How concession dynamics unfold turn-by-turn when controlling for early-settlement survivorship bias using Last-Observation-Carried-Forward (LOCF).
    \item The psychological mechanism governing the discount via statistical mediation.
    \item The paradoxical influence of baseline human-AI trust on bargaining success.
\end{enumerate}

\section{Hypotheses}
\begin{itemize}
    \item \textbf{Hypothesis 1 (Primary Price Effect):} Buyers utilizing empathy-based appeals will achieve a statistically significant lower final settlement price than buyers utilizing economic arguments ($H_1: \mu_{\text{Empathy}} < \mu_{\text{Economic}}$).
    \item \textbf{Hypothesis 2 (Buyer Surplus):} Empathy buyers will capture significantly greater economic surplus relative to their maximum private valuation ($H_2: \text{Surplus}_{\text{Empathy}} > \text{Surplus}_{\text{Economic}}$).
    \item \textbf{Hypothesis 3 (Turn Trajectory):} Under Last-Observation-Carried-Forward (LOCF) modeling, empathy appeals will induce deeper concessions across every individual round of negotiation.
    \item \textbf{Hypothesis 4 (Psychological Mediation):} The relationship between experimental condition and price concessions will be significantly mediated by perceived agent warmth ($H_4: \text{Condition} \to \text{Warmth} \to \text{Price}$).
    \item \textbf{Hypothesis 5 (AI Trust Influence):} Baseline human trust in AI systems will significantly moderate conversational engagement and price outcomes.
\end{itemize}

\section{Methodology}

\subsection{Experimental Design and Participants}
A two-condition between-subjects experimental design was deployed. Participants were randomly assigned to either the \textbf{Empathy Condition} (instructed to emphasize personal vulnerability, student/intern stipend limitations, and relationship goodwill) or the \textbf{Economic Condition} (instructed to cite local OLX market comparables, hardware age, and battery degradation). 

The experimental corpus comprises $N = 101$ total participants ($N = 99$ completed sessions: 50 Empathy, 49 Economic; 88 reached agreement, 11 impasses/walk-aways, and 2 uncompleted dropouts). All analyzed participants passed embedded pre- and post-survey attention checks ($100\%$ pass rate).

\subsection{Bargaining Task and Agent Architecture}
Participants negotiated the purchase of a used Dell Inspiron 15 laptop listed at an initial asking price of \text{₹}30,000. The buyer's private valuation was fixed at \text{₹}27,000. 

The seller agent, ``Max'', was executed via the Groq API utilizing \texttt{openai/gpt-oss-120b} at temperature $T = 0.4$. Max operated under private system instructions defining a hidden walk-away reservation floor of \text{₹}23,000 and early-turn reservation bands (prohibiting concessions below \text{₹}27,500 during Turns 1--3). Programs enforced strict programmatic anti-undercutting validation rules: any generated response that proposed an ask lower than what the buyer had already offered to pay was flagged and regenerated, preventing the LLM from bidding against itself.

\subsection{Measures and Psychometrics}
\begin{itemize}
    \item \textbf{Empathic Concern (IRI-EC):} Davis's (1980, 1983) 7-item Interpersonal Reactivity Index measuring dispositional emotional empathy. Negative items (items 2, 4, 5) were reverse-scored ($8 - x$). Scale reliability was excellent: Cronbach's $\alpha = 0.93$.
    \item \textbf{Human-AI Trust Scale:} Jian et al.'s (2000) 5-item human-machine trust scale. Item 5 was reverse-scored. Scale reliability was solid: Cronbach's $\alpha = 0.76$.
    \item \textbf{Subjective Value Inventory (Curhan's SVI):} Curhan et al.'s (2006) 8-item post-negotiation instrument measuring 4 dimensions (Instrumental Outcome, Feelings About the Self, Negotiation Process, and Relationship). Item 2 was reverse-scored. Scale reliability was near-perfect: Cronbach's $\alpha = 0.92$.
    \item \textbf{Dual Manipulation Checks:} Post-survey items assessing self-reported reliance on emotional stories versus market facts (1--7 Likert scale), supplemented by blinded keyword transcript coding of buyer messages.
\end{itemize}

\section{Results}

\subsection{Data Hygiene and Manipulation Verification}
Zero missing data or calculation errors were present in the completed dataset. Manipulation checks verified extreme compliance:
\begin{itemize}
    \item Empathy participants reported using significantly more emotional appeals ($M = 6.36, SD = 0.61$) than Economic participants ($M = 1.27, SD = 0.38$; $t(95.8) = 38.21, p < 0.001$).
    \item Economic participants reported using significantly more market/factual arguments ($M = 6.41, SD = 0.52$) than Empathy participants ($M = 1.60, SD = 0.58$; $t(95.8) = -24.80, p < 0.001$).
    \item Blinded transcript keyword coding in \texttt{code\_transcripts.js} confirmed a $98\%$ to $100\%$ text-level behavioral compliance rate.
\end{itemize}

\subsection{Primary Hypothesis: Final Deal Price}
As reported in Table~\ref{tab:primary_outcomes}, buyers in the Empathy condition negotiated a significantly lower final agreed price ($M = \text{₹}25,534.56, SD = 988.27$) compared to the Economic condition ($M = \text{₹}26,059.73, SD = 884.13$). 

A Welch's two-sample $t$-test confirmed statistical significance: $t(83.9) = -2.623, p = 0.0087$, yielding a medium-to-large effect size (Cohen's $d = -0.561$). Non-parametric Mann--Whitney testing confirmed the result without distributional assumptions ($U = 562.5, p = 0.0001$). \textbf{Hypothesis 1 is fully supported.}

\begin{table}[H]
\centering
\small
\caption{Primary and Secondary Experimental Outcomes Across Conditions}
\label{tab:primary_outcomes}
\begin{tabular}{lccccc}
\toprule
\textbf{Outcome Variable} & \textbf{Empathy ($N=43$)} & \textbf{Economic ($N=45$)} & \textbf{Test Statistic} & \textbf{$p$-value} & \textbf{Effect Size} \\
\midrule
Final Deal Price (\text{₹}) & 25,534.56 (988.27) & 26,059.73 (884.13) & $t(83.9) = -2.623$ & 0.0087** & $d = -0.561$ \\
Non-Parametric Price Rank   & Mean Rank = 37.08  & Mean Rank = 51.61  & $U = 562.5$        & 0.0001***& $r_{\text{rb}} = -0.42$ \\
Buyer Surplus (\text{₹})    & 1,465.44 (988.27)  & 940.27 (884.13)    & $t(83.9) = +2.623$ & 0.0087** & $d = +0.561$ \\
Concession Rate (\%)        & 14.88\% (3.29)     & 13.13\% (2.95)     & $t(83.9) = +2.623$ & 0.0087** & $d = +0.561$ \\
Negotiation Duration (Turns)& 5.50 (1.45)        & 5.14 (1.19)        & $t(94.2) = +1.343$ & 0.1790   & $d = +0.276$ \\
Deal Completion Rate        & 86.0\% (43/50)     & 91.8\% (45/49)     & $\chi^2(1) = 0.816$& 0.3660   & $\phi = 0.091$ \\
\bottomrule
\end{tabular}
\end{table}

\subsection{Secondary Economic Outcomes: Surplus and Intention-to-Treat}
Empathy buyers captured an average of \textbf{\text{₹}1,465.44} in economic surplus relative to their \text{₹}27,000 valuation, representing an advantage of \text{₹}525.17 over Economic buyers ($M = \text{₹}940.27, p = 0.0087$). Under an Intention-to-Treat (ITT) framework (including all 11 walk-away sessions evaluated at the seller's final ask), the empathy discount remained statistically robust: $M = \text{₹}25,809.72$ vs. \text{₹}$26,222.67$ (Mann--Whitney $p = 0.0041$).

\subsection{Concession Trajectory: Turn-by-Turn LOCF Dynamics}
Evaluating asking price trajectories across rounds using raw averages introduces severe survivorship bias: because empathy buyers settle deals early (Turns 3--4), only difficult, unyielding sessions remain active in Turns 5--6. To eliminate this distortion, we implemented \textbf{Last-Observation-Carried-Forward (LOCF)} analysis, carrying forward the settled agreement price for all subsequent rounds.

As shown in Table~\ref{tab:trajectory}, Max conceded deeper to Empathy in \textbf{every single turn} from Turn 1 through Turn 6, with the maximum divergence appearing at Turns 4--5.

\begin{table}[H]
\centering
\small
\caption{Turn-by-Turn Mean Seller Asking Price (LOCF Trajectory in \text{₹})}
\label{tab:trajectory}
\begin{tabular}{lcccc}
\toprule
\textbf{Negotiation Round} & \textbf{Empathy ($M \pm SD$)} & \textbf{Economic ($M \pm SD$)} & \textbf{Difference ($\Delta$)} & \textbf{Trajectory Pattern} \\
\midrule
Turn 0 (Listing Baseline)  & \text{₹}30,000 (0)       & \text{₹}30,000 (0)       & \text{₹}0    & Parity \\
Turn 1                     & \text{₹}28,610 (518)     & \text{₹}28,633 (476)     & -\text{₹}23  & Empathy lower \\
Turn 2                     & \text{₹}27,850 (859)     & \text{₹}27,973 (564)     & -\text{₹}123 & Empathy lower \\
Turn 3                     & \text{₹}27,064 (1,234)   & \text{₹}27,298 (681)     & -\text{₹}234 & Empathy lower \\
Turn 4                     & \text{₹}26,202 (1,137)   & \text{₹}26,615 (840)     & -\text{₹}413 & Empathy lower \\
Turn 5                     & \text{₹}25,926 (1,065)   & \text{₹}26,392 (814)     & -\text{₹}466 & Empathy lower \\
Turn 6                     & \text{₹}25,845 (1,067)   & \text{₹}26,267 (871)     & -\text{₹}422 & Empathy lower \\
\bottomrule
\end{tabular}
\end{table}

\subsection{Psychological Mediation Analysis}
To identify the mechanism through which empathy lowers prices, a Baron and Kenny path analysis was conducted testing Perceived Warmth as a mediator:
\begin{enumerate}
    \item \textbf{Path a (Condition $\to$ Warmth):} Empathy appeals dramatically elevated perceived seller warmth ($\beta = +2.407, t = 11.93, p < 0.001$).
    \item \textbf{Path b (Warmth $\to$ Price Concession):} Controlling for condition, higher perceived warmth significantly predicted lower deal prices ($\beta = -481.51, t = -2.71, p = 0.008$).
    \item \textbf{Path c (Total Effect):} The total unmediated effect of condition on final price was $\beta = -560.62$ ($p = 0.0087$).
    \item \textbf{Indirect Effect & Sobel Test:} A Sobel test confirmed significant mediation ($z = -2.48, p = 0.013$). Social warmth perception acts as the psychological vehicle unlocking LLM concessions.
\end{enumerate}

\subsection{Multiple OLS Regression and the AI Trust Paradox}
An Ordinary Least Squares (OLS) regression predicting Final Price was evaluated ($R^2 = 0.179$, Adjusted $R^2 = 0.128$; Table~\ref{tab:regression}). The model controlled for buyer age, prior negotiation experience, baseline AI trust, and initial fair price guess.

\begin{table}[H]
\centering
\small
\caption{Multiple OLS Regression: Predictors of Final Agreed Price}
\label{tab:regression}
\begin{tabular}{lccccc}
\toprule
\textbf{Predictor Variable} & \textbf{Coefficient ($\beta$)} & \textbf{Std. Error (SE)} & \textbf{$t$-statistic} & \textbf{$p$-value} & \textbf{Significance} \\
\midrule
Intercept                    & 29,160.04 & 1,940.81 & 15.025 & $< 0.001$ & $p < 0.001$*** \\
Condition (Empathy = 1)      & +35.27    & 338.95   & 0.104  & 0.9171    & Absorbed by Trust/Anchor \\
Participant Age              & +48.31    & 48.99    & 0.986  & 0.3241    & Not Significant \\
Negotiation Experience       & +138.78   & 96.61    & 1.436  & 0.1509    & Not Significant \\
\textbf{Baseline AI Trust}   & \textbf{-553.19} & \textbf{255.23} & \textbf{-2.167} & \textbf{0.0302} & \textbf{$p = 0.030$*} \\
\textbf{Fair Price Anchor}   & \textbf{-123.70} & \textbf{47.42}  & \textbf{-2.609} & \textbf{0.0091} & \textbf{$p = 0.009$**} \\
\bottomrule
\end{tabular}
\end{table}

Crucially, **Baseline Human-AI Trust** emerged as a powerful, statistically significant predictor of price discounts ($\beta = -553.19, p = 0.0302$). For every 1-point increase in a buyer's trust in AI on the Jian scale, their final deal price was \text{₹}553.19 lower.

\section{Discussion}

\subsection{The RLHF Agreeableness Bias vs. Competitive Defense}
Why does conversational empathy overpower market rationality when bargaining against an artificial counterparty? 
Two interacting mechanisms account for this finding:
\begin{enumerate}
    \item \textbf{RLHF Training Penalties:} State-of-the-art LLMs undergo extensive Reinforcement Learning from Human Feedback (RLHF) designed to align models with human values of helpfulness, harmlessness, and agreeableness. During RLHF, human evaluators heavily penalize models that appear cold, rigid, or callous when confronted with vulnerability (e.g., student exam deadlines, job layoffs, medical expenses). When buyers employ empathy, the model defaults to cooperative alignment rather than profit maximization.
    \item \textbf{Counter-Argument Activation:} When buyers utilize rational market arguments (e.g., citing cheaper OLX comparables or 85\% battery health), the LLM frames the interaction as an argumentative debate. Drawing upon its verified listing facts (Full HD screen, clean keyboard, original charger, box, and receipt), Max mounts an articulate factual defense, refusing to concede past \text{₹}26,000.
\end{enumerate}

\subsection{Unpacking the ``AI Trust Paradox''}
The finding that higher trust in AI predicts deeper price discounts ($\beta = -553.19, p = 0.0302$) resolves a key behavioral question in human-AI interaction:
\begin{itemize}
    \item \textbf{Low-Trust Cynicism and Rigid Ultimatums:} Buyers who distrust AI frequently adopt a dismissive or combative stance, viewing the machine as a simplistic algorithm that can be bullied via brute force (e.g., \textit{``You're just a bot, take ₹20,000 or leave it''}). Max's system prompt contains programmatic defense instructions to freeze asks when faced with ultimatums or unreasoned lowballs. Consequently, low-trust buyers hit deadlocks or receive minimal concessions.
    \item \textbf{High-Trust Cooperative Reciprocity:} Conversely, buyers who trust AI engage it as a legitimate bargaining partner. They negotiate in good faith, offering step-by-step concessions that align with Max's concession rules. This reciprocal engagement allows the conversation to progress smoothly to Max's lowest allowable price tiers.
\end{itemize}

\subsection{Methodological Implications: The Importance of LOCF}
Our analysis demonstrates the critical necessity of Last-Observation-Carried-Forward (LOCF) in negotiation analytics. Evaluating raw round averages created an artificial optical illusion wherein empathy appeared more expensive in Turns 5--6. This occurred purely because the most successful empathy negotiations closed at Turns 3--4, leaving only stubborn outlier sessions in later rounds. Implementing LOCF restored methodological fidelity, proving that empathy yields superior concession depth across every individual negotiation turn.

\section{Conclusion}
In human--LLM negotiation, emotional vulnerability and cooperative engagement constitute a statistically superior bargaining strategy compared to analytical economic rationality. Empathy appeals achieve a significant price discount ($p = 0.0087$), elevate buyer surplus ($p = 0.0087$), and dramatically enhance perceived social warmth ($d = 2.40$). Furthermore, human trust in AI acts not as a vulnerability, but as an asset that unlocks cooperative algorithmic reciprocity. These findings carry profound implications for the design of commercial AI agents, consumer bargaining strategies, and the governance of algorithmic marketplaces.

\end{document}
